\chapter*{KẾT LUẬN}
\addcontentsline{toc}{chapter}{KẾT LUẬN}

Báo cáo đã giúp chúng tôi hiểu rõ quy trình xây dựng mô hình dự đoán
và áp dụng một số kỹ thuật tối ưu, bao gồm tiền xử lý dữ liệu, xử lý mất cân
bằng bằng SMOTE hoặc undersampling, giảm chiều bằng PCA và huấn luyện
mô hình XGBoost. Qua đó, chúng tôi nhận thấy Accuracy không đủ để đánh
giá mô hình trong bài toán mất cân bằng; Recall, F1-score, MCC và
ROC-AUC cần được xem xét đồng thời.

Quá trình tái hiện cũng cho thấy bài báo gốc chưa mô tả đầy đủ thứ tự
chia dữ liệu, resampling, PCA và các tham số thực nghiệm. Vì vậy,
chúng tôi xây dựng một protocol minh bạch hơn: chia train/test trước,
chỉ thực hiện preprocessing và balancing trên tập huấn luyện, đồng thời
giữ nguyên phân bố lớp của tập test để tránh rò rỉ dữ liệu.

Kết quả cải thiện cho thấy việc kết hợp balancing, PCA, tuning và lựa
chọn operating threshold giúp tăng đáng kể khả năng phát hiện lớp
stroke. Trên Dataset 1, Balanced Random Forest tạo ra cải thiện nhẹ về
F1-score và MCC so với cấu hình XGBoost có trọng số. Trên Dataset 2,
XGBoost kết hợp undersampling và PCA vẫn cho sự cân bằng tốt nhất giữa
Precision, Recall, F1-score và MCC; LightGBM đạt Recall cao nhưng
Precision thấp hơn.

Từ các thí nghiệm, chúng tôi rút ra rằng không có một mô hình duy nhất
tốt nhất trên mọi bộ dữ liệu. Hiệu quả dự đoán phụ thuộc không chỉ vào
thuật toán mà còn vào protocol đánh giá, phương pháp xử lý mất cân bằng
và ngưỡng quyết định. Đây là những điểm cải thiện chính được rút ra từ
việc tái hiện và mở rộng bài báo gốc.